\section{K-029: higher odd lifts occur in the amplified component closures}

\begin{theorem}[K-029]
Let $T_k$ be the critical K-020 tensor, with
\[
 n=2+7k,\qquad D_k=2\cdot256^k,\qquad k\ge8,
 \qquad M=2^{n-1}.
\]
Let $Y_k$ be the irreducible component of $X_n^{\mu_M}$ through
its quotient point $x_k$.  This component contains a nonzero
polystable point fixed by $\mu_{3M}$.  Consequently its entire
closure is one of the collision components included in the
exceptional support of K-027 and K-028.

Writing $k=8l+s$, $0\le s<8$, the component has dimension at
least $343l^3+14l+1$, with equality when $s=0$, as in K-021.
Thus for every $l\ge1$ the actual dimension-$343l^3+14l+1$
component in tensor dimension $2+56l$ meets the higher-lift
exception.  This statement does not assert a pole at either
$M$ or $3M$.
\end{theorem}

\begin{proof}
We first give a quotient-image observation used to locate the
new point in the same component, rather than merely in the
same fixed scheme.  Suppose a finite-order product-linear
element $h$ fixes a closed tensor $T$ whose effective
special-linear stabilizer is trivial.  Let $x$ be its quotient
point, and let $\xi=\chi(h)$.  The quotient is smooth at $x$,
and its tangent representation is the quotient of the ambient
tensor tangent by the special-linear orbit tangent.  Taking
$h$-fixed vectors is exact in characteristic zero, so the
differential
\[
                         V_n^h\longrightarrow (T_xX_n)^\xi
                                                        \tag{K-029a}
\]
is surjective.  The fixed scheme is smooth at $x$.  Therefore
the restriction of the quotient map to the vector space $V_n^h$
is smooth onto that fixed germ near $T$ and its image contains
an open neighborhood of $x$ in the corresponding component.
Its image closure is irreducible and contained in $X_n^\xi$,
so it is exactly that component.  In particular, the quotient
of every vector in $V_n^h$ belongs to the same component closure.

Take $l=\lfloor k/8\rfloor\ge1$ and $s=k-8l$.  The prefix
after $k-l$ amplifier steps has dimension
\[
          m=2+7(k-l),\qquad d=7l=n-m,
 \qquad D_{k-l}=2\cdot256^{k-l}=2^sM.
                                                        \tag{K-029b}
\]
Let $e=(i,j,r)$ be its selected output edge.  The exact spectral
collapse proved in K-020--K-021 says that, after reduction modulo
$D_{k-l}$, the last $l$ gadgets replace the selected coordinate
in each factor by a block of size $1+d$.  On the original
$m$ coordinates the reduced element is an old primitive
generator, up to scalar factors with product one.  All $d$
new coordinates have the same weight as the selected old
coordinate in their respective factor.

Let $h$ be the power of the full generator of $T_k$ that induces
a primitive grading element of order $M$.  In the preceding
description it is obtained by a further $2^s$th power.  Its
restriction to the old $m$ coordinates fixes the critical prefix
tensor and has determinant-product character of order exactly
$M$.  On the new $d$ coordinates it has scalar factors
$\sigma_A,\sigma_B,\sigma_C$.  Their product is one because
the selected old edge $e$ is fixed.  Thus $h$ fixes every tensor
on the three new $d$-dimensional spaces, and its determinant
product on those spaces is
$(\sigma_A\sigma_B\sigma_C)^d=1$.

On the new spaces take the critical dimension-three tensor
$T^{(3)}$ from K-025, and add $d-3$ orthogonal one-dimensional
diagonal tensors.  Call this direct sum $Q_d$.  Its three reduced
Hermitian matrices are $I_d$.  It has an order-three fixing
element $a$ with determinant-product character of order three:
use K-025's diagonal exponents
\[
                 (0,0,1),\qquad(0,1,2),\qquad(0,1,2)
                                                        \tag{K-029c}
\]
on the first three coordinates and exponent zero on the others.
The hypothesis $d=7l\ge7$ makes this construction possible.

Now form the orthogonal direct sum
\[
                       T^*=T_{k-l}\oplus Q_d.
                                                        \tag{K-029d}
\]
All its reduced matrices are $I_n$, so its special-linear orbit
is closed and nonzero.  Both $h$ and the element acting as $a$
on the new spaces and as the identity on the old spaces fix it.
They commute, because $h$ is scalar on each new factor.  The
product has determinant-product character of order $3M$.
The quotient of $T^*$ is therefore a nonzero point of
$X_n^{\mu_{3M}}$.

On the other hand, $T^*\in V_n^h$.  The original $T_k$ is an
effectively free closed tensor by K-020, so (K-029a) places the
quotient of $T^*$ in the same component $Y_k$ through $x_k$.
This proves the collision assertion.  The dimension statement
is the exact calculation of K-021, recalled only to identify
how large these exceptional components can be.
\end{proof}

\paragraph{What this rules out.}
The high-dimensional amplified components cannot be discarded
modulo three merely because their cyclic factor-permutation
orbits have size three.  Their closures actually contain the
odd grading stabilizers that obstruct the no-collision integrality
lemma.  This is a statement about the geometry of these components;
it does not assert that their Hilbert densities have a particular
three-adic denominator or that they cancel the chain residue.
